% =====================================================================
%  Supplementary Notes (SN5): term-level standards crosswalk.
%  Referenced from the main text (sec:background) as the supplementary
%  material giving the term-level crosswalk.
% =====================================================================
\section{Standards crosswalk}
\label{sec:supp-crosswalk}

This section extends the class-level FHIR alignment in the main paper with
term-level mappings for GEM classes, properties, dimensions, and selected
enumerated values. The mappings cover the Sequence Ontology (SO), the Human
Phenotype Ontology (HPO), NCBITaxon, the Evidence and Conclusion Ontology
(ECO), SEPIO, and the GA4GH Variant Annotation (VA) model. Each mapping was
verified against its source vocabulary.

This standards crosswalk is distinct from the UMLS crosswalk in
\Cref{sec:s-umls-crosswalk}. It aligns GEM constructs with source
vocabularies, whereas that note records concept-level mappings and curation
decisions within the UMLS Metathesaurus.

The \emph{relation} column describes the relation between a GEM construct
and an external term. \emph{Exact} indicates the same scope;
\emph{close} a near match with a material scope difference;
\emph{narrower} that GEM is more specific; \emph{broader} that GEM is more
general; \emph{related} that the external term captures a constituent or
associated aspect of the GEM construct; and \emph{none} that no equivalent
has been identified. For structural correspondences that are not
term-level equivalences, the tables instead use structural labels
(\emph{structural sibling}, \emph{structural parallel}).

The crosswalk prioritizes semantic fidelity over a fixed preference for one
vocabulary. When the most faithful concept is available only in a specialized
vocabulary, the crosswalk retains that exact match rather than substituting a
broader term from a preferred vocabulary. For example,
\dimval{ANIMAL\_GENETICS} maps exactly to NCI Thesaurus
\texttt{C1510895} (\emph{Animal Genetics}), for which MeSH has no
corresponding descriptor. The \emph{relation} field records the nature of
the mapping; it does not express a preference among vocabularies.

A second source-level issue is sense merging within a single UMLS concept.
For \dimval{VARIANT}, UMLS concept \texttt{C0042333}
(\emph{Genetic Variation}) combines two senses. The MeSH descriptor defines
a population-level phenomenon (``genotypic differences observed among
individuals in a population'') and determines the concept's semantic type,
\emph{Natural Phenomenon or Process}. The NCI Thesaurus instead defines an
entity, a deviation in the nucleotide sequence of an individual, and places
the ACMG classification variants (\emph{pathogenic}, \emph{likely
pathogenic}, and \emph{benign}) directly under that concept. Its synonym
list also includes \emph{Sequence Variant}, the exact SO term shown in
\cref{tab:xwalk-dimensions}.

The UMLS mapping is therefore recorded as \emph{close}, rather than
\emph{exact}, because the concept includes the population-level MeSH sense.
We nevertheless accept its semantic type even though it lies outside the
expected axis. The entity-level NCI sense is the relevant genomics sense,
and UMLS provides no generic variant concept that is typed as an entity:
the classified child concepts are typed \emph{Nucleotide Sequence}, but
their generic parent is not.

\begin{table}[H]
  \centering
  \small
  \caption[Crosswalk: structural classes]{Structural correspondences between
  GEM and FHIR R5. The table extends the class-level alignment in the main
  paper with term-level detail.}
  \label{tab:xwalk-structural}
  \begin{tabularx}{\linewidth}{@{}>{\raggedright\arraybackslash}X >{\raggedright\arraybackslash}X l >{\raggedright\arraybackslash}X@{}}
    \toprule
    GEM construct & External term & Vocab. & Relation \\
    \midrule
        \ScientificEvidence{} (evidence item) & \texttt{Evidence} & FHIR R5 & structural sibling; GEM replaces the PICO frame with a genetic-specific dimensional frame \\
        \EvidenceVariable{} (definition and measured value) & \texttt{EvidenceVariable} + \texttt{Evidence.statistic} & FHIR R5 & broader; GEM combines the variable definition and measured value in one class \\
        Credibility (overall rating slot) & \texttt{Evidence.certainty} & FHIR R5 & structural parallel; both provide an overall rating with separable subcomponents \\
    \bottomrule
  \end{tabularx}
\end{table}

\begin{table}[H]
  \centering
  \small
  \caption[Crosswalk: assertion property and class]{Crosswalk for the
  \texttt{assertion} property and the \EvidenceAssertion{} class. The
  property records an object-level, source-anchored claim; the class is a
  meta-predicate over a decision rule.}
  \label{tab:xwalk-assertion}
  \begin{tabularx}{\linewidth}{@{}>{\raggedright\arraybackslash}X >{\raggedright\arraybackslash}X l >{\raggedright\arraybackslash}X@{}}
    \toprule
    GEM construct & External term & Vocab. & Relation \\
    \midrule
    \texttt{assertion} property (source-anchored claim) & \texttt{Statement} & GA4GH VA & narrower; a close, genetics-specific match \\
        \texttt{assertion} property & \texttt{Assertion} (\texttt{SEPIO:0000001}) & SEPIO & narrower \\
        \texttt{assertion} property (false friend) & \texttt{Evidence.assertion} (human-readable summary) & FHIR R5 & none; do not map \\
        \EvidenceAssertion{} class (meta-predicate over a decision rule) & no equivalent & --- & none \\
    \bottomrule
  \end{tabularx}
\end{table}

The \texttt{assertion} property and the \EvidenceAssertion{} class are
distinct despite their similar names. The property records the object-level,
source-anchored claim made by a paper. Its closest structural analogues are
the GA4GH VA \texttt{Statement} and the SEPIO \texttt{Assertion}, not
FHIR's \texttt{Evidence.assertion}, which is a human-readable summary.

The \EvidenceAssertion{} class is a meta-predicate over a decision rule and
is the model's point of contact with the reasoning layer
of~\cite{our_preprint}. We found no close analogue for this decision-rule
meta-predicate in FHIR, GA4GH VA, ECO, or SEPIO. Its lack of an external
analogue is therefore expected rather than a coverage failure. A future
schema revision will rename the property to remove the naming collision; the
pilot annotations will remain unchanged.

\begin{table}[H]
  \centering
  \footnotesize
  \caption[Crosswalk: dimensions and values]{GEM dimensions and selected
  values mapped to external vocabularies. Identifiers were verified against
  their source vocabularies. SO = Sequence Ontology; VA = GA4GH Variant
  Annotation.}
  \label{tab:xwalk-dimensions}
  \begin{tabularx}{\linewidth}{@{}>{\raggedright\arraybackslash}X >{\raggedright\arraybackslash}X l >{\raggedright\arraybackslash}X@{}}
    \toprule
    GEM dimension / value & External term + ID & Vocab. & Relation \\
    \midrule
        Target Type: gene & gene \texttt{SO:0000704} & SO & exact \\
        Target Type: related gene & gene \texttt{SO:0000704} (relation recorded in Gene Relation) & SO & exact \\
        Target Type: variant & sequence\_variant \texttt{SO:0001060} & SO & exact \\
        Target Type: transcript & transcript \texttt{SO:0000673} & SO & exact \\
        Target Type: interval / segment & region \texttt{SO:0000001} & SO & narrower \\
        Mode of Inheritance (axis) & Mode of inheritance \texttt{HP:0000005} & HPO & exact \\
        \quad autosomal dominant & \texttt{HP:0000006} & HPO & exact \\
        \quad autosomal recessive & \texttt{HP:0000007} & HPO & exact \\
        Organism: human & \textit{Homo sapiens} \texttt{NCBITaxon:9606} & NCBITaxon & exact \\
        Organism: mouse & \textit{Mus musculus} \texttt{NCBITaxon:10090} & NCBITaxon & exact \\
        Organism: rat & \textit{Rattus norvegicus} \texttt{NCBITaxon:10116} & NCBITaxon & exact \\
        Organism: zebrafish & \textit{Danio rerio} \texttt{NCBITaxon:7955} & NCBITaxon & exact \\
        Credibility (rating and facets) & \texttt{confidence level} (\texttt{SEPIO:0000187}), via \texttt{has\_confidence\_level} (\texttt{SEPIO:0000167}), on a rated \texttt{Assertion} (\texttt{SEPIO:0000001}) & SEPIO & exact (rating) / narrower (facets) \\
        Method: in vivo (animal model) & animal model system study evidence \texttt{ECO:0000179} & ECO & exact \\
        Method: in vitro & in vitro assay evidence \texttt{ECO:0000181} & ECO & exact \\
        Method: clinical & clinical study evidence \texttt{ECO:0000180} & ECO & exact \\
        Method: statistical genetics & no clean ECO leaf; broadest parent is \texttt{computational evidence} (\texttt{ECO:0007672}); OBI is a candidate vocabulary for study designs (unconfirmed) & ECO / OBI & none (ECO) / broader (OBI) \\
    \bottomrule
  \end{tabularx}
\end{table}

The following table lists dimensions for which the selective review reported
here identified no single external equivalent. These entries are documented
coverage gaps in the present crosswalk, not claims that no relevant term
exists in any vocabulary.

\begin{table}[H]
  \centering
  \small
  \caption[Crosswalk: coverage gaps]{Coverage gaps in the present crosswalk.
  The relation is \emph{none} in every row.}
  \label{tab:xwalk-gaps}
  \begin{tabularx}{\linewidth}{@{}l >{\raggedright\arraybackslash}X@{}}
    \toprule
    GEM dimension & Why no external term \\
    \midrule
    Knowledge Domain & no single external vocabulary represents this axis \\
        Gene Relation (conditional) & GEM-specific gene--gene relations, such as \texttt{X\_regulates\_Y} and \texttt{X\_has\_same\_function\_as\_Y}; RO may cover some relations (unconfirmed) \\
        Phenotype Scale (ordinal) & no external vocabulary provides this ordinal axis; GO biological\_process is a possible future alignment for the molecular and cellular end of the scale \\
        Variant Ascertainment & no single external term captures this dimension \\
        Resolution & represented implicitly by the selected SO target-type term, rather than as an independent external axis \\
        \EvidenceAssertion{} (meta-predicate) & belongs to a reasoning layer beyond the scope of ECO, SEPIO, VA, and FHIR \\
    \bottomrule
  \end{tabularx}
\end{table}

Taken together, the mappings and gaps show two broad patterns. First, the
dimensions that describe what the evidence concerns and which organism it
concerns map directly to established ontologies. The three experimental
Method values also map cleanly. Statistical genetics is the main Method gap:
ECO represents evidence types rather than study designs such as genome-wide
association studies and linkage analyses, and OBI is a candidate vocabulary
for that level of description.

Second, this crosswalk identifies no single external equivalent for Knowledge
Domain, Phenotype Scale, Variant Ascertainment, Resolution, or the Gene
Relation enumeration. These are GEM-specific epistemological axes, rather
than failed attempts to duplicate existing vocabularies. Their partial or
absent coverage in the reviewed standards constitutes a coverage gap in the
present crosswalk and helps motivate a dedicated reference data model.
